\scanpage{89}
et que $\inserted{G\mapsto}S\Pi M_G$ est un foncteur en \deleted{$M_G$} $G$ pour les morphismes quelconques, et que (124) est fonctoriel pour les iso. Pour un morphisme quelconque $G\xrightarrow{f}G'$, correspondant à une partie $C\subset A_G$ (telle que $G\unsure{-|\to}G'$ soit isom., en tant que quotient de $G$, à $G\unsure{-|\to}G/C$), on suppose que l'on a un diagramme commutatif
\[
\tag{125}
\diagram{\begin{tikzcd}[ampersand replacement=\&, column sep=large,row sep=large]
S\Pi M_G \arrow[r,"S\Pi(f)"] \arrow[d,"\varphi_G"] \& S\Pi M_{G'} \arrow[d,"\varphi_{G'}"] \\
S_{\unsure{C|A}}M_G \arrow[r,"\gamma_f"] \& \Pi M_{G'}
\end{tikzcd}}
\]
i.e. que le composé
\[
\diagram{\begin{tikzcd}[ampersand replacement=\&, column sep=large]
S\Pi M_G \arrow[r,"S\Pi(f)"] \& S\Pi M_{G'} \arrow[r,dash,"\varphi_{G'}"] \& \Pi M_{G'}
\end{tikzcd}}
\]
se factorise par l'extension $S_{\unsure{C|A}}M_G$ de $\Pi M_G$ par $\mathbf Z^C$, déduite de l'extension $S\Pi M_G$ de $\Pi M_G$ par $\mathbf Z^{A(G)}$ grâce à l'hom canonique
\[
\diagram{\begin{tikzcd}[ampersand replacement=\&]
\mathbf Z^{A(G)} \arrow[r] \& \mathbf Z^C.
\end{tikzcd}}
\]
On a \unsure{automatiquement|automatique} alors la relation de transitivité pour un composé.
\margin{88}
